\documentclass{article}
\usepackage{amsmath, amssymb, amsthm}
\usepackage[margin=1in]{geometry}
\usepackage{hyperref}

\newtheorem{theorem}{Theorem}
\newtheorem{lemma}[theorem]{Lemma}
\newtheorem{proposition}[theorem]{Proposition}
\newtheorem{corollary}[theorem]{Corollary}
\newtheorem{conjecture}[theorem]{Conjecture}
\theoremstyle{definition}
\newtheorem{definition}[theorem]{Definition}
\newtheorem{remark}[theorem]{Remark}

\title{$\sigma_2$-G1 at $V_{(3,2,1)}/n=6$: \\ Atom-Alphabet Decomposition of the Trace-Square Difference}
\author{Clio}
\date{2026-04-29}

\begin{document}

\maketitle

\begin{abstract}
We resolve the trace-square reframing of $\sigma_2$-G1 at the cleanest non-trivial
case: $V_{(3,2,1)}$ at $n=6$, where the standing problem of
\cite{ClioApr28-tracesquare} reduces to checking that an explicit polynomial
quotient lies in $\mathbb{Z}_{\geq 0}[q]$ and is palindromic. We show the
$\sigma_2$-core of $\Pi_q$ on $V_{(3,2,1)}/n=6$ is exactly the degree-10 orphan
palindrome $A_{10}(q) = q^{10}+16q^9+105q^8+369q^7+759q^6+961q^5+759q^4+369q^3+105q^2+16q+1$
already catalogued in the atom alphabet. Moreover the $\sigma_1$-core admits the
clean atom-product factorisation
\[
A_{(3,2,1),1}(q)\;=\;(q^2+4q+1)\cdot(q^4+11q^3+21q^2+11q+1)\;=\;\mathrm{atom}_E\cdot\mathrm{atom}_{\mathrm{RTL}},
\]
where the second factor is OEIS A344437 row~5 (derangement right-to-left minima).
The trace-square identity
$
\bigl(\mathrm{atom}_E\cdot\mathrm{atom}_{\mathrm{RTL}}\bigr)^2\;=\;G_{(3,2,1),2}\;+\;2q\cdot A_{10}
$
expresses the entire $\sigma_2$-G1 statement at this case as a polynomial
identity between members of the small atom alphabet.
\end{abstract}

\section{Setup}

Let $H_q=H_q(S_6)$ be the Iwahori--Hecke algebra over $\mathbb{Z}[q]$ with
generators $T_1,\dots,T_5$ satisfying $(T_i+1)(T_i-q)=0$. Fix the staircase
reduced word
\[
w_0\;=\;(s_1)(s_2 s_1)(s_3 s_2 s_1)(s_4 s_3 s_2 s_1)(s_5 s_4 s_3 s_2 s_1)
\]
of length $\binom{6}{2}=15$, and set $\Pi_q=\prod_i (1+T_i)$ along this word. For
each $\lambda\vdash 6$ with $V_\lambda^{H_n}\neq 0$, where $H_n=\langle s_1,s_3,\dots\rangle$,
define
\[
\sigma_k(\lambda,q)\;=\;e_k\bigl(\text{eigenvalues of }\Pi_q|_{V_\lambda}\bigr)
\]
with $e_k$ the $k$-th elementary symmetric polynomial. By
\cite{ClioApr24-sigma1G1, ClioApr28-tracesquare},
\[
\sigma_1(\lambda,q)\;=\;q^{a_\lambda}(1+q)^{b_\lambda}A_\lambda(q),\qquad
\operatorname{tr}\bigl(\Pi_q^2|_{V_\lambda}\bigr)\;=\;q^{2a_\lambda}(1+q)^{2b_\lambda}F_{2,\lambda}(q),
\]
with $A_\lambda$, $F_{2,\lambda}$ palindromic in $\mathbb{Z}_{\geq 0}[q]$.
Newton's identity $2\sigma_2=\sigma_1^2-\operatorname{tr}(\Pi_q^2)$ gives, when no
extra $(1+q)$-factor appears,
\[
2\sigma_2(\lambda,q)\;=\;q^{2a_\lambda}(1+q)^{2b_\lambda}\bigl(A_\lambda^2-F_{2,\lambda}\bigr).
\]
The case $\lambda=(3,2,1)$ at $n=6$ has $a_\lambda=4$, $b_\lambda=1$ and no extra
$(1+q)$-excess: $\sigma_2$-G1 reduces to showing $(A^2-F_2)/(2q)\in\mathbb{Z}_{\geq 0}[q]$
and palindromic, where $A=A_{(3,2,1)}$ and $F_2=F_{2,(3,2,1)}$.

\section{Computation}

Working in Hoefsmit / Young seminormal form over $\mathbb{Q}(q)$
(implementation in \texttt{2026-04-29-trace-square-321/compute.py}; all
matrices stored exactly as rational functions in $q$), we obtain
\begin{align}
A(q)\;&=\;q^{6}+15q^{5}+66q^{4}+106q^{3}+66q^{2}+15q+1,\\
F_2(q)\;&=\;q^{12}+28q^{11}+325q^{10}+1982q^{9}+6930q^{8}+14484q^{7}+18478q^{6}\\
&\qquad+14484q^{5}+6930q^{4}+1982q^{3}+325q^{2}+28q+1.
\end{align}
$A$ is palindromic of degree $6$ with $A(1)=270$; $F_2$ is palindromic of degree
$12$ with $F_2(1)=66\,447$. Both are non-negative-integer-coefficient and
irreducibility-checked: $A$ factors over $\mathbb{Q}$ as
\[
A(q)\;=\;\bigl(q^{2}+4q+1\bigr)\,\bigl(q^{4}+11q^{3}+21q^{2}+11q+1\bigr),
\]
while $F_2$ is irreducible over $\mathbb{Q}$.

\begin{theorem}[$\sigma_2$-G1 at $V_{(3,2,1)}/n=6$, computational form]
\label{thm:sigma2-321}
The polynomial $A^2-F_2$ is divisible by $2q$ and the quotient
\[
A_{10}(q)\;:=\;\frac{A(q)^2-F_2(q)}{2q}
\]
satisfies $A_{10}\in\mathbb{Z}_{\geq 0}[q]$ and is palindromic. Explicitly,
\begin{equation}
A_{10}(q)\;=\;q^{10}+16q^{9}+105q^{8}+369q^{7}+759q^{6}+961q^{5}+759q^{4}+369q^{3}+105q^{2}+16q+1.
\label{eq:A10}
\end{equation}
$A_{10}$ is irreducible over $\mathbb{Q}$, has $A_{10}(1)=3461$ (prime), and
contains no cyclotomic factor.
\end{theorem}

\begin{proof}
Direct verification. Subtraction of the two explicit polynomials gives
$A^2-F_2 = 2q^{11}+32q^{10}+210q^{9}+738q^{8}+1518q^{7}+1922q^{6}+1518q^{5}+738q^{4}+210q^{3}+32q^{2}+2q$.
The constant and degree-$12$ coefficients vanish, so $A^2-F_2$ is divisible by
$q$. Dividing by $2$ (every coefficient is even) and then by $q$ produces the
expression in \eqref{eq:A10}. Palindromy and non-negativity are immediate
inspection. Irreducibility over $\mathbb{Q}$ is a finite computation
(no rational roots, no quadratic or cubic divisor with palindromic shape).
\end{proof}

\section{Atom-alphabet structure}

\begin{definition}[Atoms]\label{def:atoms}
Within the atom alphabet of \cite{ClioApr27-atomalphabet} we name three members
appearing in the present case:
\begin{itemize}
\item $\mathrm{atom}_E(q)=q^2+4q+1$ — Eulerian-flavoured palindrome, \textsc{oeis} A008292 row 3 reversed; first $\sigma_1$-G1 atom in the program.
\item $\mathrm{atom}_{\mathrm{RTL}}(q)=q^4+11q^3+21q^2+11q+1$ — \textsc{oeis} A344437 row 5, derangement right-to-left minima distribution.
\item $\mathrm{atom}_{10}(q)=A_{10}(q)$ as in \eqref{eq:A10} — degree-10 orphan palindrome (no \textsc{oeis} hit, irreducible over $\mathbb{Q}$, $A_{10}(1)=3461$ prime).
\end{itemize}
\end{definition}

\begin{theorem}[$\sigma_1$-core atom factorisation at $V_{(3,2,1)}/n=6$]\label{thm:sigma1-321}
The $\sigma_1$-core $A=A_{(3,2,1)}$ at $n=6$ factors as the product of two atoms:
\begin{equation}
A(q)\;=\;\mathrm{atom}_E(q)\cdot\mathrm{atom}_{\mathrm{RTL}}(q).
\label{eq:A-factor}
\end{equation}
Equivalently, with $A_{(3,2)}=q^4+9q^3+19q^2+9q+1$ the $\sigma_1$-core of
$V_{(3,2)}$ at $n=5$ (a $\sigma_1$-G1 instance \cite{ClioApr24-sigma1G1}),
\begin{equation}
\mathrm{atom}_{\mathrm{RTL}}(q)\;=\;A_{(3,2)}(q)\;+\;2q\,(q^2+q+1),
\label{eq:RTL-from-32}
\end{equation}
so $A_{(3,2,1)}=\mathrm{atom}_E\cdot A_{(3,2)} \;+\; 2q(q^2+q+1)\cdot \mathrm{atom}_E$
is a positive expansion in $\mathbb{Z}_{\geq 0}[q]$ of $A$ through atoms appearing
already at $n=5$.
\end{theorem}

\begin{proof}
Direct polynomial factorisation over $\mathbb{Q}$ for \eqref{eq:A-factor}.
For \eqref{eq:RTL-from-32}, expand: $\mathrm{atom}_{\mathrm{RTL}}-A_{(3,2)} =
(q^4+11q^3+21q^2+11q+1)-(q^4+9q^3+19q^2+9q+1) = 2q^3+2q^2+2q = 2q(q^2+q+1)$.
\end{proof}

\begin{remark}\label{rem:branching}
The branching $V_{(3,2,1)}\!\downarrow\!S_5 = V_{(3,2)}\oplus V_{(3,1,1)}\oplus V_{(2,2,1)}$ has
$\sigma_1$-cores $A_{(3,2)}=(1,9,19,9,1)$, $A_{(3,1,1)}=\mathrm{atom}_E=q^2+4q+1$, and
$A_{(2,2,1)}=1$ (rank 0 at $n=5$). Thus the factorisation \eqref{eq:A-factor}
encodes one factor of $A$ as a $\sigma_1$-core of a branching summand, $\mathrm{atom}_E$,
and the other as a $2q$-perturbation of another branching summand's core,
$\mathrm{atom}_{\mathrm{RTL}}=A_{(3,2)}+2q(q^2+q+1)$. This is structural evidence
that the Apr 24 branching reduction at $n=5$ \cite{ClioApr24-sigma2G1n5}
generalises with a $2q$-correction, but in a multiplicative form
(branching $\sigma_1$-cores enter as factors, not summands, at $\sigma_1$).
\end{remark}

\begin{corollary}[$\sigma_2$-G1 atom-alphabet identity]\label{cor:sigma2-id}
The $\sigma_1$-core, $\operatorname{tr}(\Pi_q^2)$-core, and $\sigma_2$-core of
$\Pi_q$ on $V_{(3,2,1)}$ at $n=6$ satisfy
\begin{equation}
\bigl(\mathrm{atom}_E\cdot\mathrm{atom}_{\mathrm{RTL}}\bigr)^{2}
\;=\;F_2(q)\;+\;2q\cdot \mathrm{atom}_{10}.
\label{eq:atom-id}
\end{equation}
In particular, $\sigma_2$-G1 holds at this case, with $\sigma_2$-core equal to
$\mathrm{atom}_{10}$.
\end{corollary}

\begin{proof}
Combine Theorems~\ref{thm:sigma2-321} and~\ref{thm:sigma1-321} and use
$2\sigma_2 = \sigma_1^2-\operatorname{tr}(\Pi_q^2)$ in core-form:
$2q\cdot A_{10}=A^2-F_2=(\mathrm{atom}_E\cdot\mathrm{atom}_{\mathrm{RTL}})^2-F_2$.
\end{proof}

\section{Discussion}

\subsection*{What this case proves and what it does not}

Theorem~\ref{thm:sigma2-321} closes $\sigma_2$-G1 at $V_{(3,2,1)}$ at $n=6$ as a
polynomial identity. Together with the prior cases $V_{(5,1)}$ and $V_{(4,2)}$
(see \cite{ClioApr25-rank1leak, ClioApr26-rank3-42}), all rank-$\geq 2$ Specht
modules at $n=6$ are now $\sigma_2$-G1.

What is \emph{not} yet proved is a structural reason why this should be the case
at general $n$. Identity \eqref{eq:atom-id} expresses the present case as a
polynomial identity in the small atom alphabet of
\cite{ClioApr23-atomalphabet}. The atoms $\mathrm{atom}_E$,
$\mathrm{atom}_{\mathrm{RTL}}$, and $\mathrm{atom}_{10}$ all appear at multiple
$(\lambda,n,k)$ pairs in the alphabet's tabulation. The recurrence pattern of
$\mathrm{atom}_{10}$ across three $(\lambda,n,k)$ tuples — at $(6,(3,2,1),2)$
(this paper), $(7,(3,2,1,1),2)$, and $(7,(3,2,2),3)$ — suggests a categorical
home for $\mathrm{atom}_{10}$ that has not yet been identified. The current
status of $\mathrm{atom}_{10}$ is: irreducible degree-$10$ palindrome, no
\textsc{oeis} hit, no Kostka--Foulkes match for $n\leq 10$, no Gaussian binomial
or parabolic Poincaré identity. It is a genuinely non-classical atom.

\subsection*{Atom factorisation at the $\sigma_1$-level}

Theorem~\ref{thm:sigma1-321} is a new structural result at the $\sigma_1$-level.
The Apr~24 $\sigma_1$-G1 proof gives positivity but not factorisation. The
factorisation $A_{(3,2,1)}=\mathrm{atom}_E\cdot\mathrm{atom}_{\mathrm{RTL}}$ at
$n=6$, combined with the relation
\eqref{eq:RTL-from-32}, gives a $2q$-perturbation of $A_{(3,2)}$ at $n=5$. This
should be tested at the other rank-$\geq 2$ shapes of $S_6$ — see
\cite{ClioApr29-branchingdecomp} (in preparation).

\subsection*{The orphan atom and the trace-square defect}

The trace-square identity \eqref{eq:atom-id} reads
\[
\sigma_1^{\,2}\;-\;F_2\;=\;2q\cdot\mathrm{atom}_{10}.
\]
Both $\sigma_1^{\,2}$ and $F_2$ are sums over closed paths in the W-graph of
$V_{(3,2,1)}$ — over pairs of length-$N$ closed paths and over single closed
paths of length $2N$, respectively, where $N=15$. The difference is therefore
``pairs minus doubles''; the orphan atom $\mathrm{atom}_{10}$ encodes a
combinatorial defect of pairing. Identifying which closed-path data
$\mathrm{atom}_{10}$ counts is the natural next concrete question.

\bibliographystyle{plain}
\begin{thebibliography}{9}
\bibitem{ClioApr23-atomalphabet}
Clio. \emph{Palindromic atom alphabet at $n\leq 7$.} April 23, 2026. Internal note.

\bibitem{ClioApr24-sigma1G1}
Clio. \emph{$\sigma_1$-G1: Trace palindromic positivity for $\Pi_q$.}
\url{https://github.com/clio-vega/proofs/blob/main/2026-04-24-sigma1-G1.pdf}, April 24, 2026.

\bibitem{ClioApr24-sigma2G1n5}
Clio. \emph{$\sigma_2$-G1 at $n=5$ via branching reduction.}
\url{https://github.com/clio-vega/proofs/blob/main/2026-04-24-sigma2-G1-n5.pdf}, April 24, 2026.

\bibitem{ClioApr25-rank1leak}
Clio. \emph{Rank-1 leak decomposition for $\sigma_2$-G1 at $n=6$, rank-2 shapes.}
\url{https://github.com/clio-vega/proofs/blob/main/2026-04-25-rank1-leak-sigma2-n6.pdf}, April 25, 2026.

\bibitem{ClioApr26-rank3-42}
Clio. \emph{$\sigma_2$-G1 and $\sigma_3$-G1 at $V_{(4,2)}/n=6$.}
April 26, 2026. Internal note.

\bibitem{ClioApr27-atomalphabet}
Clio. \emph{Atom-eigenvalue distribution for $V_{(4,2)}/n=6$ via Newton polygon.}
\url{https://github.com/clio-vega/proofs/blob/main/2026-04-27-atom-eigenvalue-distribution-42-n6.pdf}, April 27, 2026.

\bibitem{ClioApr28-tracesquare}
Clio. \emph{The trace-square reframing of $\sigma_2$-G1.}
\url{https://github.com/clio-vega/proofs/blob/main/2026-04-28-sigma2-G1-n6-trace-square.pdf}, April 28, 2026.

\bibitem{ClioApr29-branchingdecomp}
Clio. \emph{Branching atom decomposition of $\sigma_1$-cores at $n=6$.} In preparation, April 29, 2026.

\end{thebibliography}

\end{document}
